\section{Finite witness amalgamation by viability}
\label{sec:amalgamation}

\subsection{Layered finite relations}

Fix a numerical parameter \(p\) and a horizon \(h\).  At level \(k\le h\),
answers are numbers \(a<B(p,k)\).  Two Boolean predicates are given:
\[
 \Good(p,k,a),\qquad
 \Compat(p,k,a,b),
\]
where the second predicate is used only when
\(a<B(p,k)\) and \(b<B(p,k+1)\).  A path of length \(d\) beginning at level
\(k\) is a tuple
\[
 (a_0,\ldots,a_d),
 \qquad a_j<B(p,k+j),
\]
such that every \(a_j\) is good and every adjacent pair is compatible.
Tuples are coded by any fixed primitive-recursive list or bounded-sequence
coding.  The exact coding is immaterial once projection, length, and bounded
folds have their defining equations.

A na\"ive attempt to construct a path chooses a good answer at each level and
hopes that the choices remain extendible.  This can fail even in a finite
binary tree: a locally good node may be a dead end.  Backward viability removes
this failure without using any choice principle.

\subsection{Backward viability}

Define \(\Viable(p,k,d,a)\) recursively in the remaining depth \(d\):
\begin{align*}
 \Viable(p,k,0,a)
   &:= [a<B(p,k)]\wedge \Good(p,k,a),\\
 \Viable(p,k,d+1,a)
   &:= [a<B(p,k)]\wedge \Good(p,k,a)\wedge{}\\
   &\quad \Any_{b<B(p,k+1)}
      \bigl(\Good(p,k+1,b)\wedge
             \Compat(p,k,a,b)\wedge
             \Viable(p,k+1,d,b)\bigr).
\end{align*}
Every connective and bounded existential in this definition is a Boolean fold,
so \(\Viable\) is primitive recursive.  The changing level causes no problem:
the recursion variable is \(d\), while \(k\) is an ordinary parameter.

\begin{theorem}[Finite path--viability equivalence]
\label{thm:path-viability}
For every fixed layered presentation, iPRA proves
\begin{align*}
 \Viable(p,k,d,a)=1
 \quad\Longleftrightarrow\quad&
 \text{there is a coded compatible good path of length }d\\
 &\text{beginning with }a\text{ at level }k.
\end{align*}
The witness path may be bounded by a primitive-recursive code bound depending
only on \(p,k,d\) and the candidate bounds \(B\).
\end{theorem}

\begin{proof}
Induct on \(d\).  At depth zero, the singleton path \([a]\) is valid exactly
when \(a\) is in range and good.  For the successor direction, unfold the
recursion equation for \(\Viable\).  Correctness of the finite \(\Any\)-fold
gives a witness \(b<B(p,k+1)\) satisfying the edge condition and
\(\Viable(p,k+1,d,b)=1\).  The induction hypothesis yields a path beginning
with \(b\); prepend \(a\).  Conversely, a path beginning with \(a,b\) supplies
the required bounded witness \(b\), while its tail supplies viability by the
induction hypothesis.

The induction formula can be Booleanized as equality of two checkers: the
viability bit and the bounded path-existence bit.  Hence the formal proof uses
quantifier-free induction.  The displayed existential statement is recovered
from the search-correctness theorem after the Boolean induction is complete.
\end{proof}

\subsection{A canonical finite path}

Let
\[
 \Next(p,k,d,a)
\defeq \Search_{b<B(p,k+1)}
  \bigl(\Good(p,k+1,b)\wedge\Compat(p,k,a,b)
       \wedge\Viable(p,k+1,d,b)\bigr).
\]
When \(\Viable(p,k,d+1,a)=1\), bounded-search correctness proves that
\(\Next(p,k,d,a)\) is a compatible viable successor.  Primitive recursion on
\(j\le d\), repeatedly applying \(\Next\), gives a canonical path
\(\mathsf{path}(p,k,d,a,j)\).

\begin{corollary}[Least viable path selector]
\label{cor:least-viable-path}
There is a fixed primitive-recursive path constructor such that iPRA proves
\[
 \Viable(p,k,d,a)=1\to
 \mathsf{PathOK}(p,k,d,a,\mathsf{path}(p,k,d,a))=1.
\]
No uniqueness hypothesis is needed.  The least-witness convention is only a
finite implementation choice; the mathematical conclusion remains relational.
\end{corollary}

\begin{proof}
Use quantifier-free induction on the number of generated edges.  The invariant
states that the current node is viable for the number of edges still required
and that the accumulated path is good and compatible.  Search correctness
establishes the successor invariant.  At the final stage, depth-zero viability
supplies goodness of the last node.  All path-prefix assertions are finite
Boolean folds.
\end{proof}

\subsection{From local seriality to finite chains}

Sometimes the input theorem supplies not a precomputed viable root but a
seriality contract:
\begin{align*}
 &\exists a<B(p,0)\;\Good(p,0,a),\tag{I}\label{eq:serial-init}\\
 &\forall k<h\;\forall a<B(p,k)\;
   \bigl(\Good(p,k,a)\to
     \exists b<B(p,k+1)\,[\Good(p,k+1,b)\wedge\Compat(p,k,a,b)]\bigr).
   \tag{E}\label{eq:serial-step}
\end{align*}
For a fixed finite horizon, bounded search converts these assumptions to a
path without countable choice.

\begin{theorem}[Finite dependent choice is primitive recursive]
\label{thm:finite-dependent-choice}
For every horizon \(h\), assumptions \eqref{eq:serial-init} and
\eqref{eq:serial-step} yield a primitive-recursively constructed compatible
good path through levels \(0,\ldots,h\).  The theorem is provable in iPRA.
\end{theorem}

\begin{proof}
Choose the least initial good candidate by bounded search.  At stage \(k<h\),
specialize the seriality premise to the current candidate and choose the least
compatible successor by bounded search.  Induct on \(k\) with the
quantifier-free invariant that the coded prefix has the required length and
its Boolean path checker equals one.  The horizon is a free numerical
parameter; there is no infinite sequence and no appeal to countable choice.
\end{proof}

Backward viability is stronger than this serial construction.  It works when
some good nodes are dead ends, and it supports additional global constraints
that can be tested only after later levels are known.

\subsection{Constraint propagation and witness amalgamation}

Let \(C(p,k,a)\) be an additional Boolean constraint, perhaps a prescribed
residual budget, chart membership, sign bracket, or compatibility with a
second proof module.  Replace \(\Good\) by \(\Good\wedge C\) in the viability
recursion.  This gives the following monotonicity and amalgamation principle.

\begin{theorem}[Finite witness amalgamation]
\label{thm:witness-amalgamation}
Suppose a finite collection of requirements is compiled into one layered
finite relation.  If the initial viability fold is nonempty,
\[
 \Any_{a<B(p,0)}\Viable(p,0,h,a)=1,
\]
then bounded search and the least viable path constructor produce one finite
witness family satisfying all requirements simultaneously.  Adding a new
Boolean requirement can only remove viable nodes; it cannot create a false
path.
\end{theorem}

\begin{proof}
Search supplies a viable initial node.  Apply
\cref{cor:least-viable-path}.  Soundness of conjunction and of every folded
requirement shows that the output path satisfies each component.  For
monotonicity, an induction on remaining depth proves that viability under the
stronger goodness predicate implies viability under the weaker one.  At the
successor step, the same witness edge is used.
\end{proof}

This is a finite constraint-satisfaction theorem.  It applies not only to
linear paths but, after coding the state of a bounded collection of pending
requirements, to finite trees, finite acyclic proof networks, overlapping
chart choices, and simultaneous approximation windows.  The state-space may
be large, but it remains primitive recursive and its correctness proof is the
same Boolean invariant.

\subsection{Full-horizon viability}

The next section needs to know not merely that a node extends forward, but that
it occurs on a complete path from level zero to a specified horizon.  Define
\[
 \FullV(p,h,k,a)=1
\]
to mean that \(k\le h\) and there is a compatible good path
\((a_0,\ldots,a_h)\) with \(a_k=a\).  This is a bounded path-existence
predicate and hence primitive recursive.  It also admits a dynamic-programming
factorization:
\[
 \FullV(p,h,k,a)
 \quad\Longleftrightarrow\quad
 \Reach(p,k,a)\wedge\Viable(p,k,h-k,a),
\]
where \(\Reach\) is the forward Boolean fold saying that \(a\) is reachable
from some good level-zero node.  Both directions are proved by concatenating
or splitting finite paths.

This two-sided notion is crucial.  A merely forward-viable candidate at a high
level need not be connected to the input data.  A merely reachable candidate
may be a dead end.  Full-horizon viability records exactly the candidates that
participate in a complete finite solution.

\subsection{Object-language implementation}

There are two efficient implementations, both suitable for the trace compiler.

\paragraph{Path-code checker.}
Bound a path code, search over all codes, and verify range, goodness,
compatibility, and the distinguished coordinate by finite folds.  This version
has the smallest correctness proof and is useful for certificates.

\paragraph{Dynamic-programming table.}
Build the forward reachable table and backward viable table, each represented
as a bounded bit-vector code.  The local invariant states that the first
\(j\) entries agree with the defining fold.  This version avoids exponential
search and is appropriate for compiled algorithms.

In either implementation, the mathematical theorem is an instance of
\cref{thm:embedded-trace}.  The application module must supply object-language
derivations for Boolean lookup, update, folds, and the local recurrence; the
generic trace theorem then supplies the global soundness derivation.
